\documentclass{article}
\usepackage{geometry}
\geometry{letterpaper, margin=0.9in, includefoot}

\usepackage{forest}
\usepackage{lmodern}
\usepackage{fontspec}
\setmainfont{Calibri}
\setmonofont{Consolas}

\usepackage{amsmath}
\usepackage{booktabs}
\usepackage{xcolor}
\usepackage{listings}
\usepackage{scrextend}
\changefontsizes{14.4pt}
\setlength{\emergencystretch}{5em} % Prevent overfull lines
\setlength{\hbadness}{10000} % Allow some overfull lines

\usepackage{microtype}
\usepackage[colorlinks=true, linkcolor=blue, urlcolor=blue]{hyperref}

% Custom listing styling for project structure and code blocks
% Define custom colors for syntax highlighting
\definecolor{codegray}{rgb}{0.83,0.83,0.83}
\definecolor{codeblue}{rgb}{0.1,0.1,0.8}
\definecolor{codegreen}{rgb}{0.0,0.5,0.0}
\definecolor{codepurple}{rgb}{0.58,0,0.82}
\definecolor{backcolour}{rgb}{0.95,0.95,0.92}
\definecolor{codeborder}{rgb}{0.85,0.85,0.88}

\lstdefinestyle{codeblock}{
    backgroundcolor=\color{codegray},
    basicstyle=\ttfamily\small,
    breaklines=true,
    keepspaces=true,
    frame=single,
    rulecolor=\color{codeborder},
    framerule=0.5pt,
    framesep=8pt,
    tabsize=2,
    showstringspaces=false
}

% Define JavaScript Language for listings
\lstdefinelanguage{JavaScript}{
  keywords={break, case, catch, continue, debugger, default, delete, do, else, finally, for, function, if, in, instanceof, new, return, switch, this, throw, try, typeof, let, const, var, void, while, with, async, await, export, import, from},
  morecomment=[l]{//},
  morecomment=[s]{/*}{*/},
  morestring=[b]',
  morestring=[b]",
  morestring=[b]`,
  sensitive=true
}

% Configure global styling for code blocks
\lstset{
  language=JavaScript,
  backgroundcolor=\color{backcolour},   
  commentstyle=\color{codegreen},
  keywordstyle=\color{codeblue},
  numberstyle=\tiny\color{codegray},
  stringstyle=\color{codepurple},
  basicstyle=\ttfamily\footnotesize,
  breakatwhitespace=false,         
  breaklines=true,                 
  captionpos=b,                    
  keepspaces=true,                 
  numbers=left,                    
  numbersep=5.5pt,                  
  showspaces=false,                
  showstringspaces=false,
  showtabs=false,                  
  tabsize=2,
  frame=single,
  rulecolor=\color{black!30}
}

\title{\textbf{Personal Portfolio Website (\texttt{rah3.github.io})}}
\author{\textbf{Rahie Datta}}
\date{Carleton University - Bachelor of Computer Science}

\begin{document}

\maketitle

\noindent\textbf{Live Site:} \url{https://rah3.github.io} \\
\textbf{GitHub Repository:} \url{https://github.com/rah3/rah3.github.io}

\vspace{0.5cm}

\section{About}

I am a fourth-year Computer Science student at \textbf{Carleton University} pursuing a Bachelor of Computer Science with a Minor in Mathematics. 

Core technical interests and focus areas include:
\begin{itemize}
    \item Software Engineering \& System Design
    \item Linux Systems \& Process Management
    \item Backend Development \& APIs
    \item Database Design \& Optimization
    \item Telecommunications Systems
    \item Automation \& Scripting
\end{itemize}

\section{Technology Stack}

The project relies on a modern, lightweight, and framework-free JavaScript architecture. It uses Node.js and Express solely for local development ergonomics, while the production assets are fully optimized for static delivery via GitHub Pages.

\subsection{Frontend}
\begin{itemize}
    \item \textbf{HTML5 / CSS3}: Semantic markup and modern custom styling.
    \item \textbf{Vanilla JavaScript (ES6+)}: Modular, performant implementation without heavy framework overhead.
\end{itemize}

\subsection{Backend \& Tooling}
\begin{itemize}
    \item \textbf{Node.js \& Express}: Local development server and static asset proxying.
    \item \textbf{Git \& GitHub}: Version control and CI/CD hosting via GitHub Pages.
    \item \textbf{Nodemon}: Hot-reloading utility during local development.
\end{itemize}

\section{Architectural Style}

The codebase implements a \textbf{modular component pattern} and functional composition. Rather than relying on heavy client-side frameworks, the UI components are isolated into independent, testable modules.

\subsection{Core Modules}
\begin{itemize}
    \item \texttt{Navigation.js}: Manages mobile responsive toggle states and layout adjustments.
    \item \texttt{ScrollSpy.js}: Tracks user scroll position to dynamically highlight active navigation links.
    \item \texttt{EmailClipboard.js}: Handles user interactions for quick email copying with visual feedback.
\end{itemize}

\subsection{Application Factory Pattern}
Application startup and lifecycle management are initialized through a lightweight factory pattern:

\begin{lstlisting}[style=codeblock, language=JavaScript]
createApp().init();
\end{lstlisting}

\section{Features}
Here are the features that make this portfolio website stand out:
\begin{itemize}
    \item \textbf{Fully Responsive}: Mobile-first design adapting seamlessly across mobile, tablet, and desktop viewports.
    \item \textbf{Active Section Highlighting}: Real-time viewport tracking using intersection observers / scroll spy.
    \item \textbf{Interactive Utilities}: One-click email copy-to-clipboard functionality.
    \item \textbf{Lightweight \& Fast}: Zero frontend framework bloat ensuring instant page load speeds.
    \item \textbf{Express Development Server}: Local mock environment mirroring production static behavior.
\end{itemize}

\section{Project Structure}
Below is a visual representation of the repository structure:
\begin{center}
\begin{forest}
  for tree={
    font=\ttfamily\scriptsize,
    grow'=0,
    child anchor=west,
    parent anchor=south,
    anchor=west,
    calign=first,
    l sep=15pt,
    s sep=6pt,
  }
  [rah3.github.io/
    [{index.html -- Main entry markup}]
    [{style.css -- Global stylesheet}]
    [js/
      [{App.js -- Core application factory}]
      [{Navigation.js -- Mobile menu controller}]
      [{ScrollSpy.js -- Active section tracker}]
      [{EmailClipboard.js -- Clipboard utility}]
    ]
    [{server.js -- Express local development server}]
    [{package.json -- Node.js dependencies and scripts}]
    [{README.md -- Project documentation}]
  ]
\end{forest}
\end{center}

\section{Featured Projects}
Below are some of the notable projects I have worked on during my academic career:
\begin{enumerate}
    \item Telecommunications Switching Database (COMP 3005)
    \item Linux Systems \& Shell Development (COMP 3000)
    \item Professional Portfolio Website (rah3.github.io)
\end{enumerate}

\subsection{Telecommunications Switching Database}
\begin{itemize}
    \item \textbf{Course:} COMP 3005 (Database Management Systems)
    \item \textbf{Technologies:} SQLite, SQL, Relational Database Design
    \item \textbf{Overview:} Designed and implemented a robust relational database schema tailored for telecommunications routing and subscriber switching management.
\end{itemize}

\subsection{Linux Systems \& Shell Development}
\begin{itemize}
    \item \textbf{Course:} COMP 3000 (Operating Systems)
    \item \textbf{Technologies:} Linux, Bash, System Calls, C
    \item \textbf{Overview:} Explored low-level OS primitives, process management (\texttt{fork}, \texttt{execve}), environment variable handling, and remote SSH administration.
\end{itemize}

\subsection{Professional Portfolio Website}
\begin{itemize}
    \item \textbf{Technologies:} Node.js, Express, JavaScript, HTML5/CSS3
    \item \textbf{Overview:} This repository was built from scratch to highlight academic and practical engineering work.
\end{itemize}

\section{Local Development}

To run the development environment locally:

\begin{enumerate}
    \item \textbf{Clone the repository:}
\begin{lstlisting}[style=codeblock]
git clone https://github.com/rah3/rah3.github.io.git
cd rah3.github.io
\end{lstlisting}

    \item \textbf{Install dependencies:}
\begin{lstlisting}[style=codeblock]
npm install
\end{lstlisting}

    \item \textbf{Start the development server:}
\begin{lstlisting}[style=codeblock]
npm run dev
\end{lstlisting}
    \textit{(Or standard start: \texttt{npm start})}

    \item \textbf{Open in your browser:} \\
    Navigate to \url{http://localhost:3000}.
\end{enumerate}

\section{Deployment}

This repository uses a decoupled deployment model:
\begin{itemize}
    \item \textbf{Production:} The contents of the \texttt{public/} directory are deployed statically to \textbf{GitHub Pages}.
    \item \textbf{Development:} The \texttt{server.js} Express backend is used exclusively for local testing and does not run in production.
\end{itemize}

\section{Author}

\textbf{Rahie Datta} \\
Carleton University - Bachelor of Computer Science (BCS) \\
Expected Graduation: 2027 \\
GitHub: \url{https://github.com/rah3}

\end{document}